
\section{Related Work}

\todo{Strategies with Tracing JIT, PyPy}
\todo{Strategies with Method JIT (Truffle, V8)}
\todo{Pointer tagging}

\subsection{Graal and Truffle}

The Graal/Truffle project is a VM building framework similar to RPython in its main design goal.
An interpreter written in a statically typed programming language is used as the source artifact to automatically produce an efficient VM including a JIT.

Graal \cite{graal} is an experimental extension to the HotSpot VM which exposes VM internal mechanisms to the executed program via a Java API.
Namely, HotSpots method based JIT compiler can be controlled from the client program.
Truffle \cite{truffle} is a generic, self-optimizing AST interpreter that makes use of the interface provided by Graal.
In order to implement a VM using Graal/Truffle, the program code must be parsed into a syntax tree and passed to Truffle for execution.
Truffle will dynamically modify the AST using techniques like partial evaluation and inlining to create compilation units for heavily used code paths, which are then converted to machine code by the Graal JIT.
This approach effectively uses a method JIT compiler, but is similar to RPythons meta tracing JIT.

The main difference between RPython and Graal/Truffle, aside from the underlying infrastructure, lies in the flexibility for the VM programmer.
Graal/Truffle requires the VM implementor to use an AST, while the RPython toolchain accepts arbitrary Python programs as long as they satisfy the static typechecking.

The project JRuby+Truffle \footnote{\url{http://www.chrisseaton.com/rubytruffle/}} is a Ruby VM written in Java which uses the Graal/Truffle infrastructure.
JRuby+Truffle includes storage strategy optimizations \cite{smallDataStructures}, but not in form of a reusable library like rstrategies.

